\documentclass[11pt]{article}
\usepackage{amsmath, amssymb, geometry, hyperref}
\geometry{margin=1in}

\title{\textbf{Thermal Rigidity: Linking Metric Saturation to the Zero-Kelvin Limit and Quantum Superposition}}
\author{Research Manuscript Supplement}
\date{2026}

\begin{document}
\maketitle

\begin{abstract}
This paper explores the correspondence between the mechanical saturation of the vacuum and the thermodynamic limit of absolute zero. We propose that the kinetic arrest observed at 0 K is a manifestation of absolute metric rigidity ($P_{max}$), providing a structural explanation for the stability of quantum superposition.
\end{abstract}

\section{The Thermodynamic Manifestation of $P_{max}$}
We propose that temperature $T$ is functionally linked to the elasticity of the Tensor Field ($\xi$). As a localized region approaches maximum saturation, the field's capacity to yield to constituent vibration diminishes:
\begin{equation}
\lim_{\xi \to 0} T = 0 \text{ K}
\end{equation}
In the "Frozen Core" of a high-density singularity or a micro-frozen state, the internal kinetic energy is arrested by the infinite rigidity of the medium.



\section{Mechanical Basis for Superposition}
Superposition requires a reduction in environmental "noise" or decoherence. We suggest that cryogenic cooling (achieving 0 K) is an artificial simulation of the "Frozen Core" state. By reducing constituent vibration, we minimize the field's kinetic resets, allowing a constituent to exist in multiple displacement states. This suggests that the center of a black hole may exist in a state of macroscopic perfect superposition due to absolute metric saturation.

\section{Disclosure and Sentiment}
This framework is presented as an indicative study into the material properties of the vacuum. The author suggests that the "Quantum Limit" may be fundamentally linked to a "Mechanical Limit" of the underlying tensor field. No conflict of interest is declared.

\begin{thebibliography}{9}
\bibitem{Unruh1976} Unruh, W. G. (1976). Notes on black-hole evaporation. \textit{Phys. Rev. D}.
\bibitem{Zurek2003} Zurek, W. H. (2003). Decoherence, einselection, and the quantum origins of the classical. \textit{Rev. Mod. Phys.}
\end{thebibliography}

\end{document}